\section{Dirichlet 特征与同余子群的 Eisenstein 级数}
\label{sec:eisenstein-congruence}
% ============================================================================

到目前为止，我们讨论的模形式都住在最对称、也最理想化的世界里：全模群 $\SL_2(\Z)$。
但真正的算术问题——尤其是和椭圆曲线、Galois 表示以及费马大定理相关的问题——往往迫使我们走向更精细的群，例如 $\GammaO(N)$、$\GammaI(N)$。
一旦群变小，"允许的对称性"也变得更复杂，我们就不能再指望所有 Eisenstein 级数都像全模群情形那样只用一个公式解决。
新的问题是：\textbf{当我们对格点施加同余条件时，应该用什么语言把这些条件编码进去？}

答案就是 Dirichlet 特征。
你可以先把它想成一种"会记住模 $N$ 余数类的乘性权重"：整数不再只是被当成一个数，而是连同它模 $N$ 的行为一起被记录下来。
这正适合处理同余子群，因为很多限制本来就是"$d\equiv 1\pmod N$"、"$c\equiv 0\pmod N$"这类条件；而角色分解的核心工具——特征的正交性——恰好能把这种同余限制拆成一串可操作的加权和。

学会 Dirichlet 特征以后，我们立刻能做三件以前做不到的事。
第一，我们可以系统地构造 $\Gamma_0(N)$、$\Gamma_1(N)$ 上的 Eisenstein 级数，而不只是停留在全模群。
第二，这些扭曲 Eisenstein 级数的 $q$-系数会自然变成扭曲除数和
$\sum_{d\mid n}\chi_1(d)\chi_2(n/d)d^{k-1}$，让同余信息直接进入 Fourier 系数。
第三，Dirichlet 特征本身也定义了 $L(s,\chi)$，于是 Eisenstein 级数与 Dirichlet $L$-函数这两条主线就在这里汇合。
从这个角度看，本节不是额外加一个定义，而是在为"模形式与 $L$-函数如何对话"铺路。

对全模群 $\SL_2(\Z)$，我们已经完全理解了 Eisenstein 级数 $E_k$。
但本书的最终目标（费马大定理）需要考虑同余子群，
特别是 $\GammaO(N)$、$\GammaI(N)$。
这些子群上的 Eisenstein 级数由 \textbf{Dirichlet 特征}（Dirichlet character）参数化。

\begin{definition}[Dirichlet 特征]
\label{def:dirichlet-character}
\index{Dirichlet 特征 (Dirichlet character)}\index{chi@$\chi$}
设 $N$ 为正整数。\textbf{模 $N$ 的 Dirichlet 特征}是函数 $\chi\colon \Z \to \C$，满足：
\begin{enumerate}
  \item \textbf{周期性}：$\chi(n + N) = \chi(n)$ 对所有 $n$；
  \item \textbf{完全乘性}：$\chi(mn) = \chi(m)\chi(n)$ 对所有 $m, n$；
  \item \textbf{与 $N$ 互素时非零}：$\chi(n) \neq 0 \iff \gcd(n, N) = 1$；
    当 $\gcd(n, N) > 1$ 时 $\chi(n) = 0$。
\end{enumerate}
特征 $\chi$ 本质上是群同态 $(\Z/N\Z)^\times \to \C^\times$，
延拓到 $\Z$ 上（不与 $N$ 互素的地方补零）。

模 $N$ 的特征共有 $\varphi(N) = |(\Z/N\Z)^\times|$ 个。
其中 $\chi_0$（满足 $\chi_0(n) = 1$ 对所有 $\gcd(n,N)=1$）称为\textbf{主特征}。
\end{definition}

% figures/ch04-dirichlet-character.tex
% Dirichlet 特征 χ mod 5 的值示意表
\begin{figure}[ht]
\centering
\begin{tikzpicture}[font=\small, >=Stealth]

% === 标题 ===
\node[font=\normalsize\bfseries] at (0, 2.8) {$(\Z/5\Z)^\times$ 上的特征};

% 单位根在单位圆上的位置
\draw[gray!40] (0,0) circle (1.5cm);
\draw[->] (-2, 0) -- (2, 0) node[right, font=\footnotesize] {$\Re$};
\draw[->] (0, -2) -- (0, 2) node[above, font=\footnotesize] {$\Im$};

% 四个四次单位根
\foreach \ang/\val/\lab in {
  0/{$1$}/{$\chi(1)=1$},
  90/{$i$}/{$\chi(2)=i$},
  180/{$-1$}/{$\chi(3)=-1$},
  270/{$-i$}/{$\chi(4)=-i$}
}{
  \fill[blue!70!black] ({1.5*cos(\ang)}, {1.5*sin(\ang)}) circle (3pt);
  \node[font=\footnotesize] at ({1.9*cos(\ang)}, {1.9*sin(\ang)}) {\val};
}

% 对应的 a mod 5
\node[red!70!black, font=\footnotesize] at (1.5+0.4, -0.3) {$a=1$};
\node[red!70!black, font=\footnotesize] at (0.2, 1.5+0.3) {$a=2$};
\node[red!70!black, font=\footnotesize] at (-1.5-0.5, -0.3) {$a=3$};
\node[red!70!black, font=\footnotesize] at (0.2, -1.5-0.3) {$a=4$};

% 生成元 2 的标注
\draw[->, thick, orange!70!black, bend left=20]
  ({1.5*cos(0)+0.1}, {1.5*sin(0)}) to ({1.5*cos(90)}, {1.5*sin(90)+0.1});
\node[orange!70!black, font=\footnotesize] at (1.2, 0.9) {$\times 2$};

% === 右侧：表格 ===
\begin{scope}[xshift=5cm, yshift=0.3cm]
\node[font=\footnotesize\bfseries] at (0, 2.2) {特征值表};

% 表格内容
\foreach \row/\a/\ca/\cb/\cc/\cd in {
  0/{$a \bmod 5$}/{$1$}/{$2$}/{$3$}/{$4$},
  -0.7/{$\chi_0$ (主特征)}/{$1$}/{$1$}/{$1$}/{$1$},
  -1.4/{$\chi_1$}/{$1$}/{$i$}/{$-i$}/{$-1$},
  -2.1/{$\chi_2$}/{$1$}/{$-1$}/{$-1$}/{$1$},
  -2.8/{$\bar\chi_1$}/{$1$}/{$-i$}/{$i$}/{$-1$}
}{
  \node[font=\footnotesize, anchor=east] at (-1.2, \row) {\ca};
  \node[font=\footnotesize] at (-0.4, \row) {\cb};
  \node[font=\footnotesize] at (0.4, \row) {\cc};
  \node[font=\footnotesize] at (1.2, \row) {\cd};
  \node[font=\footnotesize, anchor=east] at (-2.0, \row) {\a};
}
% 分隔线
\draw[gray] (-2.3, -0.35) -- (1.5, -0.35);
\draw[gray] (-2.3, 1.8) -- (1.5, 1.8);

% 列标签
\node[font=\footnotesize\bfseries, anchor=east] at (-2.0, 1.9) {特征 $\backslash$ $a$};

\end{scope}

\end{tikzpicture}
\caption{模 $5$ 的 Dirichlet 特征。$(\Z/5\Z)^\times = \{1,2,3,4\}$
是循环群，生成元 $2$（因为 $2^1=2, 2^2=4, 2^3=3, 2^4=1 \pmod 5$）。
共有 $4$ 个特征，对应四次单位根 $\{1, i, -1, -i\}$（左侧单位圆）。
$\chi_0$ 是主特征，$\chi_1$ 是将生成元 $2$ 映到 $i$ 的特征。
所有特征延拓到 $\Z$ 上时，令 $\chi(n) = 0$ 若 $\gcd(n,5)>1$。}
\label{fig:dirichlet-character}
\end{figure}


\begin{example}[模 $4$ 的特征]
\label{ex:chi-mod4}
$(\Z/4\Z)^\times = \{1, 3\}$，是二阶群。恰好有两个特征：
\begin{itemize}
  \item 主特征 $\chi_0$：$\chi_0(1) = \chi_0(3) = 1$，$\chi_0(2) = \chi_0(0) = 0$。
  \item 非主特征 $\chi_{-1}$（奇特征）：$\chi_{-1}(1) = 1$，$\chi_{-1}(3) = -1$，
    $\chi_{-1}(2) = 0$。
\end{itemize}
对整数 $n$：$\chi_{-1}(n) = \left(\frac{-4}{n}\right)$（Kronecker 符号），
即奇数时为 $\pm 1$（$n \equiv 1 \pmod 4$ 时为 $+1$，$n \equiv 3$ 时为 $-1$），偶数时为 $0$。
\end{example}

\begin{example}[特征的正交性]
\label{ex:character-orthogonality}
\index{正交关系 (orthogonality relations)}
模 $N$ 的特征满足\textbf{正交关系}：
\[
  \frac{1}{\varphi(N)} \sum_{a \in (\Z/N\Z)^\times} \chi(a) \overline{\psi(a)}
  = \begin{cases} 1 & \text{若 } \chi = \psi, \\ 0 & \text{若 } \chi \neq \psi. \end{cases}
\]
以及对固定 $a \in (\Z/N\Z)^\times$，
\[
  \frac{1}{\varphi(N)} \sum_{\chi \bmod N} \chi(a) \overline{\chi(b)}
  = \begin{cases} 1 & \text{若 } a \equiv b \pmod N, \\ 0 & \text{否则}. \end{cases}
\]
这是有限 Abel 群上 Fourier 分析的特例，是 Dirichlet $L$-函数理论的基础。
\end{example}

\begin{definition}[Dirichlet $L$-函数]
\label{def:dirichlet-L-function}
\index{Dirichlet L-函数 (Dirichlet L-function)}\index{Ls@$L(s,\chi)$}
设 $\chi$ 是模 $N$ 的 Dirichlet 特征。对 $\mathrm{Re}(s) > 1$，定义
\[
  L(s, \chi) = \sum_{n=1}^{\infty} \frac{\chi(n)}{n^s}
  = \prod_{p \nmid N} \frac{1}{1 - \chi(p)p^{-s}} \cdot \prod_{p \mid N} \frac{1}{1 - \chi(p)p^{-s}}.
\]
（后者是 \textbf{Euler 乘积}，来自 $\chi$ 的完全乘性。）

主特征时：$L(s, \chi_0) = \zeta(s) \prod_{p \mid N}(1 - p^{-s})$（从 $\zeta(s)$ 去掉有限个因子）。

$L(s, \chi)$ 可以亚纯延拓到全复平面，当 $\chi \neq \chi_0$ 时是整函数，
满足函数方程（类比 $\zeta(s)$ 的函数方程）。
\end{definition}

\begin{example}[$L(1, \chi_{-1})$ 与 $\pi$]
\label{ex:L-function-value}
取 $\chi = \chi_{-1}$（模 $4$ 的奇特征）：
\[
  L(1, \chi_{-1}) = 1 - \frac{1}{3} + \frac{1}{5} - \frac{1}{7} + \cdots = \frac{\pi}{4}.
\]
这是 Leibniz 公式！Dirichlet $L$-函数在 $s = 1$ 处的值与算术有深刻联系：
\[
  L(1, \chi) \neq 0 \quad \text{对所有非主特征 } \chi.
\]
这正是 \textbf{Dirichlet 素数定理}（$\gcd(a,N)=1$ 的等差数列 $a, a+N, a+2N, \ldots$
中有无穷多个素数）的关键引理。
\end{example}

\textbf{同余子群上的 Eisenstein 级数。}
现在用 Dirichlet 特征构造 $\GammaO(N)$ 和 $\GammaI(N)$ 的 Eisenstein 级数。

设 $\chi_1, \chi_2$ 分别是模 $M_1, M_2$ 的原始特征
（即不被更小模的特征"诱导"得到），$N = M_1 M_2$。
当 $\chi_1(-1)\chi_2(-1) = (-1)^k$ 时（奇偶相容条件），定义

\begin{definition}[扭曲 Eisenstein 级数]
\label{def:twisted-eisenstein}
\index{扭曲 Eisenstein 级数 (twisted Eisenstein series)}
\[
  E_{k,\chi_1,\chi_2}(\tau)
  = c_0 + \sum_{n=1}^{\infty} \left(\sum_{d \mid n} \chi_1(n/d)\chi_2(d)\, d^{k-1}\right) q^n,
\]
其中 $c_0 = 0$（若 $M_1 > 1$）或 $c_0 = L(1-k, \chi_2)$（若 $M_1 = 1$）。
这里 $L(1-k, \chi_2)$ 是 Dirichlet $L$-函数在非正整数处的特殊值（由解析延拓定义）。
\end{definition}

\begin{example}[最简单的扭曲——$\chi_1 = \chi_0, \chi_2 = \chi_0$（全模群）]
当 $M_1 = M_2 = 1$（两个平凡特征），
$\sum_{d \mid n} \chi_0(n/d)\chi_0(d) d^{k-1} = \sigma_{k-1}(n)$，
$c_0 = L(1-k, \chi_0) = \zeta(1-k) = -B_k/k$，
这就回到了 $G_k/(2\zeta(k)) = E_k$。
\end{example}

\begin{proposition}[扭曲 Eisenstein 级数是模形式]
\label{prop:twisted-eisenstein-modular}
设 $\chi_1(-1)\chi_2(-1) = (-1)^k$，$N = M_1 M_2$。则
$E_{k,\chi_1,\chi_2} \in \Mk_k(\GammaI(N), \chi_1\chi_2)$，
即满足
\[
  E_{k,\chi_1,\chi_2}(\gamma(\tau)) = \chi_1\chi_2(d)\,(c\tau+d)^k\, E_{k,\chi_1,\chi_2}(\tau)
  \quad \text{对} \quad \gamma = \begin{pmatrix} a & b \\ c & d \end{pmatrix} \in \GammaI(N).
\]
\end{proposition}

\begin{proof}[证明思路]
仿照 $G_k$ 的证明，把 $E_{k,\chi_1,\chi_2}$ 写成格求和：
\[
  G_{k,\chi_1,\chi_2}(\tau)
  = \sum_{\substack{c,d \in \Z \\ (c,d) \neq (0,0)}}
    \frac{\chi_1(c)\chi_2(d)}{(c\tau + d)^k}.
\]
当 $\gamma_0 = \begin{pmatrix} a & b \\ c_0 & d_0 \end{pmatrix} \in \GammaI(N)$（$d_0 \equiv 1 \pmod N$）时，
变量替换 $(c,d) \to (ca + dc_0, cb + dd_0)$ 把 $\chi_1(c)\chi_2(d)$ 变成
$\chi_1(c)\chi_2(d) \cdot \chi_2(d_0)$（因为 $dd_0 \equiv d \pmod{M_2}$），
给出 $\chi_2(d_0) = \chi_1\chi_2(d_0)$（由 $d_0 \equiv 1 \pmod{M_1}$）。
\end{proof}

\textbf{Eisenstein 级数的完备性。}
$\Mk_k(\Gamma)$ 有一个重要分解：
\[
  \Mk_k(\Gamma) = \mathcal{E}_k(\Gamma) \oplus \Sk_k(\Gamma),
\]
其中 $\mathcal{E}_k(\Gamma)$（\textbf{Eisenstein 子空间}）由上述扭曲 Eisenstein 级数张成，
$\Sk_k(\Gamma)$（\textbf{尖点形式子空间}）是它的正交补（在 Petersson 内积下，见第三章）。
这个分解说明 Eisenstein 级数"捕获"了模形式空间的"非尖点"部分。

\subsection*{$L$-函数特殊值、类数公式与 BSD 的原型}

CSS 的一个重要风格，是把 Dirichlet $L$-函数的特殊值看成后面更高阶现象的原型。对次数 $1$ 的 Artin 对象，特殊值给出类数；对次数 $2$ 的模形式和椭圆曲线，Birch--Swinnerton-Dyer 猜想则预言特殊值给出秩、调和子群与 Tate--Shafarevich 群等更精细的不变量。

\begin{definition}[广义 Bernoulli 数]
\label{def:generalized-bernoulli}
设 $\chi$ 是导子为 $f$ 的 Dirichlet 特征。定义广义 Bernoulli 数 $B_{m,\chi}$ 使得
\[
  \sum_{a=1}^{f}\chi(a)\frac{te^{at}}{e^{ft}-1}
  = \sum_{m=0}^{\infty} B_{m,\chi}\frac{t^m}{m!}.
\]
它们把通常的 Bernoulli 数（对应平凡特征）推广到了带特征的情形。
\end{definition}

\begin{proposition}[Dirichlet $L$-函数的一般特殊值公式]
\label{prop:dirichlet-special-values}
设 $\chi$ 是导子为 $f$ 的原始特征，$m\ge 1$。则
\[
  L(1-m,\chi)=-\frac{B_{m,\chi}}{m}.
\]
若再记 Gauss 和
\[
  \tau(\chi)=\sum_{a=1}^{f}\chi(a)e^{2\pi ia/f},
\]
则函数方程把正整数点也化成显式公式：当 $\chi(-1)=(-1)^m$ 时，
\[
  L(m,\chi)
  =(-1)^{1+\frac{m-\epsilon}{2}}\frac{\tau(\chi)}{2}
   \left(\frac{2\pi i}{f}\right)^m
   \frac{B_{m,\overline{\chi}}}{m!},
  \qquad \epsilon\in\{0,1\},\ \chi(-1)=(-1)^\epsilon.
\]
因此 $L(m,\chi)$ 总是“$\pi^m$ 乘以代数数”。
\end{proposition}

\begin{proof}[说明]
第一条是 Dirichlet $L$-函数在非正整数处的标准插值公式；第二条则由完成 $L$-函数的函数方程把 $L(m,\chi)$ 反射到 $L(1-m,\overline{\chi})$ 得到。对本书后续内容而言，关键不是记住每个常数，而是要记住：\textbf{特殊值可以被完全显式地算出来。}
\end{proof}

\begin{theorem}[Dirichlet 类数公式]
\label{thm:dirichlet-class-number}
设 $D<0$ 是基本判别式，$\chi_D=\left(\frac{D}{\cdot}\right)$ 是对应的二次特征。则
\[
  L(1,\chi_D)=\frac{2\pi\,h(D)}{w(D)\sqrt{|D|}},
\]
其中 $h(D)$ 是虚二次域 $\Q(\sqrt{D})$ 的类数，$w(D)=|\mathcal{O}_D^\times|$ 是单位个数。
\end{theorem}

\begin{example}[两个最经典的特殊值]
\label{ex:quadratic-L-values}
\leavevmode
\begin{itemize}
  \item $D=-4$ 时，$h(-4)=1$，$w(-4)=4$，所以
    \[
      L(1,\chi_{-4})=\frac{2\pi}{4\cdot 2}=\frac{\pi}{4}.
    \]
    这正是 Leibniz 级数
    \[
      1-\frac13+\frac15-\frac17+\cdots=\frac{\pi}{4}.
    \]
  \item $D=-3$ 时，$h(-3)=1$，$w(-3)=6$，所以
    \[
      L(1,\chi_{-3})=\frac{2\pi}{6\sqrt3}=\frac{\pi}{3\sqrt3}.
    \]
\end{itemize}
\end{example}

\textbf{注。} 这个类数公式可以看成 Birch--Swinnerton-Dyer 猜想的“次数 $1$ 原型”：$L(1,\chi_D)$ 这样的解析量，精确等于类数 $h(D)$ 这样的代数量。到了椭圆曲线情形，BSD 仍然说“特殊值等于算术不变量”，只是右边会变得更复杂。


% ============================================================================
